\documentclass{article}

% NeurIPS 2025/2026 style (submission mode with line numbers)
\usepackage{neurips_2025}  % NeurIPS style file (using 2025 template)

% Standard packages
\usepackage[utf8]{inputenc}
\usepackage[T1]{fontenc}
\usepackage{hyperref}
\usepackage{url}
\usepackage{booktabs}
\usepackage{amsfonts}
\usepackage{amsmath}
\usepackage{amssymb}
\usepackage{nicefrac}
\usepackage{microtype}
\usepackage{xcolor}
\usepackage{graphicx}
\usepackage{pifont}  % for \ding checkmarks
\usepackage{enumitem}
\usepackage{multirow}
\usepackage{makecell}
\usepackage{subcaption}
\usepackage{algorithm}
\usepackage{algorithmic}
\usepackage{amsthm}

\newtheorem{definition}{Definition}

% Custom commands
\newcommand{\drift}{\textsc{IntentDrift}}
\newcommand{\bench}{\textsc{IntentDrift-SWE}}
\newcommand{\ie}{\textit{i.e.}}
\newcommand{\eg}{\textit{e.g.}}
\newcommand{\cf}{\textit{cf.}}
\newcommand{\etal}{\textit{et al.}}

\title{IntentDrift-SWE: Evaluating Software Engineering Agents \\ Under Dynamic Requirement Evolution}

\author{
  Anonymous Author(s)
}

\begin{document}

\maketitle

\begin{abstract}
LLM-powered software engineering agents are increasingly deployed to assist developers in fixing bugs, implementing features, and refactoring code within real repositories. Existing benchmarks such as SWE-bench evaluate these agents under an implicit assumption: the user's requirement is static and fully specified from the outset. In practice, however, requirements dynamically evolve during multi-turn coding sessions---developers progressively narrow the scope of a bug, add constraints after partial implementation, shift to an unrelated task mid-session, express contradictory goals, or harbor implicit expectations (e.g., tests, documentation) that surface only through interaction. This mismatch leads to systematic overestimation of agent capabilities.

We introduce \bench{}, a benchmark for evaluating SWE agents under dynamic requirement evolution. Our contributions are threefold. First, we propose a \emph{taxonomy of six intent drift patterns}---progressive refinement, topic shift, constraint addition, goal conflict, implicit need emergence, and intent backtracking---grounded in real developer workflows with AI coding assistants. Second, we construct an executable benchmark of 100+ tasks across 5--8 popular Python repositories, each equipped with a Docker environment, an environment-coupled user simulator that triggers drift based on the agent's code state, and deterministic test-based evaluation with the $\text{pass}^k$ reliability metric. Third, we propose complementary process metrics---Drift Detection Latency (DDL), Context Preservation Score (CPS), and Rollback Rate---all computed deterministically without LLM-as-judge. Evaluating 5 frontier models, we find that requirement drift causes a substantial performance drop compared to static baselines, with constraint addition and intent backtracking proving most challenging. Our results demonstrate that current SWE benchmarks significantly overestimate agents' real-world capability to handle the iterative, evolving nature of software development.
\end{abstract}

% ============================================================
\section{Introduction}
\label{sec:introduction}

Consider the following scenario: a developer asks an SWE agent to ``fix this bug in the login flow.'' The agent begins investigating, locates a suspicious null-pointer dereference in the session handler, and proposes a patch. The developer reviews the diff and clarifies: ``actually, the issue is in the auth module---the token validation is wrong.'' The agent pivots, now examining authentication logic. Mid-way through the fix, the developer adds a constraint: ``don't change the public API---we have downstream consumers.'' The agent revises its approach accordingly. After seeing the updated patch, the developer reconsiders: ``use the other approach you mentioned earlier---the one with the middleware wrapper.'' The agent must now backtrack, discard its current changes, and reconstruct a solution that honors all accumulated constraints.

In existing SWE benchmarks~\citep{jimenez2024swebench, yang2024sweagent}, this interaction would never arise. SWE-bench evaluates agents on isolated GitHub issues: one issue description maps to one gold patch, with no user interaction whatsoever. This single-turn, static-intent formulation has driven remarkable progress---yet it systematically exempts agents from one of the most fundamental challenges of real-world software engineering assistance: keeping up with a developer whose understanding of the problem evolves as the agent works on it. We term this phenomenon \textbf{\drift{}}---the dynamic evolution of user intent during multi-turn interactions with SWE agents.

\paragraph{The static intent assumption in SWE benchmarks.}
The dominant evaluation paradigm for SWE agents assumes that the developer's intent is fully and correctly specified at the outset. SWE-bench~\citep{jimenez2024swebench} provides a single issue description and evaluates the agent's patch against repository test suites, with no opportunity for clarification or refinement. SWE-agent~\citep{yang2024sweagent} introduces a sophisticated agent-computer interface but inherits the same single-turn evaluation protocol. More recent efforts have begun to relax certain static assumptions: SWE-EVO~\citep{sweevo2025} studies cross-version evolution of codebases, testing whether agents can adapt patches as repositories change over time, but does not model interactive user intent changes within a single task. SWE-CI~\citep{sweci2026} introduces continuous integration feedback loops, but the requirements are generated from automated CI signals rather than from an evolving human developer. Neither benchmark captures the iterative, conversational process through which developers and SWE agents collaboratively arrive at solutions.

\paragraph{Why \drift{} matters for SWE agents.}
Intent drift poses uniquely severe challenges in software engineering contexts. Unlike conversational assistants that produce text responses, SWE agents execute \emph{concrete, often irreversible actions}: they edit source files, delete code, refactor APIs, modify configurations, and run test suites whose side effects persist. When a developer's intent shifts after the agent has already committed changes---renaming functions that other modules depend on, deleting files that turn out to be needed, or restructuring code in ways that are difficult to undo---the agent faces a fundamentally harder problem than simply updating its next response. It must reason about which prior edits remain valid under the new intent, which must be rolled back, and how to reconcile the revised goal with the current repository state. A failure to handle such drift gracefully can leave the codebase in an inconsistent state, introduce subtle regressions, or silently violate constraints the developer assumed were being respected. These failure modes are invisible to benchmarks that assume static, single-turn intents.

\paragraph{Our contributions.}
We introduce \textbf{\drift{}}, a systematic framework for evaluating SWE agents under dynamic user intent evolution. Our work makes the following contributions:

\begin{itemize}[leftmargin=*, itemsep=2pt]
    \item \textbf{A taxonomy of intent drift} (\S\ref{sec:taxonomy}): We identify and formalize six distinct patterns of intent drift---progressive refinement, topic shift, constraint addition, goal conflict, implicit need emergence, and intent backtracking---grounded in HCI literature and analysis of real developer--agent interactions. Each drift event is further decomposed along five orthogonal dimensions (distance, explicitness, timing, frequency, reversibility), enabling fine-grained characterization.

    \item \textbf{\bench{}-SWE} (\S\ref{sec:benchmark}): We construct a benchmark of 120 multi-turn SWE tasks spanning 6 popular Python repositories, each annotated with drift types, drift points, per-turn intent states, and expected repository end-states verified by test suites. Tasks are constructed via a two-stage \emph{blueprint-to-trajectory} pipeline with three-layer verification (automated test validation, cross-model LLM review, and dual human annotation with Cohen's $\kappa \geq 0.75$).

    \item \textbf{Deterministic, drift-aware metrics} (\S\ref{sec:metrics}): We propose five evaluation metrics grounded in test-based, deterministic assessment: \emph{pass\textsuperscript{k}} (adapted from $\tau$-bench~\citep{yao2024taubench}), measuring the probability that all $k$ sampled trajectories pass the final test suite; \emph{Drift Detection Latency} (DDL), quantifying turns elapsed before the agent acknowledges a drift event; \emph{Context Preservation Score} (CPS), assessing retention of still-valid prior context across drift boundaries; \emph{Rollback Rate}, measuring how often the agent correctly undoes invalidated changes; and \emph{Proactive Clarification Rate} (PCR), capturing unprompted requests for disambiguation.

    \item \textbf{Systematic evaluation and findings} (\S\ref{sec:experiments}): We evaluate multiple frontier models and SWE agent frameworks, including SWE-agent~\citep{yang2024sweagent}, Claude Code, and Cursor. Key findings include: (1) intent drift causes substantial pass\textsuperscript{k} degradation compared to static-intent baselines, (2) constraint addition and intent backtracking are the most challenging drift types, as they require reasoning about action reversibility in code, (3) implicit drift is significantly harder than explicit drift across all agents, and (4) agents with richer context management (\eg, file-level undo, checkpoint mechanisms) show disproportionate advantages on backtracking scenarios.
\end{itemize}

\paragraph{Positioning.}
We distinguish \drift{} from adjacent efforts along two axes. First, unlike SWE-EVO~\citep{sweevo2025} and SWE-CI~\citep{sweci2026}, which study agent robustness to \emph{codebase-level} or \emph{CI-driven} changes, we focus on \emph{user-side} intent evolution within a single interactive session. Second, unlike $\tau$-bench~\citep{yao2024taubench} and $\tau$-Knowledge~\citep{tauknowledge2026}, which evaluate tool agents in customer-service domains with simulated users, we target the software engineering domain where actions are persistent, state-modifying, and subject to hard correctness constraints via test suites. Our benchmark and metrics are designed to expose the gap between static evaluations and the interactive reality of SWE agent deployment.

% ============================================================

% ============================================================
\section{Related Work}
\label{sec:related}

\bench{} connects four research threads: SWE agent evaluation, tool-agent-user interaction benchmarks, user intent in software engineering, and agent behavioral drift. We survey each and highlight the gap our work fills.

% ==============================================================
\subsection{SWE Agent Evaluation}
\label{sec:rw-swe}
% ==============================================================

SWE-bench~\citep{jimenez2024swebench} established the dominant paradigm for evaluating autonomous software engineering agents: given a single GitHub issue, produce a single patch, and validate it against held-out tests. The dataset comprises 2{,}294 tasks drawn from real-world pull requests across 12 Python repositories. SWE-bench Verified curates a 500-task subset with human verification by 93 annotators, improving label quality and enabling more reliable comparisons. SWE-agent~\citep{yang2024sweagent} introduced the Agent--Computer Interface (ACI) design philosophy, demonstrating that purpose-built interfaces significantly boost agent performance (NeurIPS 2024). On the data side, SWE-smith~\citep{swesmith2025} scales training-data synthesis by systematically generating bug--fix pairs, earning a NeurIPS 2025 D\&B Spotlight. To address data contamination, SWE-bench Live~\citep{swebenchlive2025} provides a continuously updated stream of post-training-cutoff issues. SWE-EVO~\citep{sweevo2025} probes robustness to cross-version evolution: across 48 tasks GPT-4o's resolve rate drops from 65\% to 21\% when the target repository version changes. Most recently, SWE-CI~\citep{sweci2026} extends evaluation to CI-loop maintenance with an Architect--Programmer dual-agent setup over 100 tasks, introducing the EvoScore metric for measuring adaptation across build cycles.

\begin{table}[t]
\centering
\small
\caption{Comparison of SWE evaluation benchmarks. \bench{} is the first to model multi-turn user interaction with evolving requirements during coding.}
\label{tab:swe-compare}
\begin{tabular}{lccccc}
\toprule
\textbf{Benchmark} & \textbf{Tasks} & \textbf{Multi-turn} & \textbf{User} & \textbf{Requirement} & \textbf{Process} \\
 & & \textbf{Interaction} & \textbf{Simulation} & \textbf{Evolution} & \textbf{Metrics} \\
\midrule
SWE-bench           & 2{,}294 & \ding{55} & \ding{55} & \ding{55} & \ding{55} \\
SWE-bench Verified  & 500     & \ding{55} & \ding{55} & \ding{55} & \ding{55} \\
SWE-bench Live      & ongoing & \ding{55} & \ding{55} & \ding{55} & \ding{55} \\
SWE-smith           & 50k+    & \ding{55} & \ding{55} & \ding{55} & \ding{55} \\
SWE-EVO             & 48      & \ding{55} & \ding{55} & Partial  & \ding{55} \\
SWE-CI              & 100     & Partial  & \ding{55} & Partial  & \ding{51} \\
\midrule
\textbf{\bench{}}   & XXX     & \ding{51} & \ding{51} & \ding{51} & \ding{51} \\
\bottomrule
\end{tabular}
\end{table}

As Table~\ref{tab:swe-compare} shows, a critical assumption pervades the entire SWE-bench lineage: \emph{requirements are given once and never change}. Every benchmark supplies a fixed issue description and measures whether the agent produces a correct patch. None models the user who filed the issue---there is no interaction, no clarification, and no possibility that the requirement evolves mid-task. Real-world software development, by contrast, is fundamentally iterative: developers routinely refine specifications, shift priorities, and add constraints after coding has begun. \bench{} fills this gap by introducing multi-turn user interaction with structured requirement evolution into SWE agent evaluation.

% ==============================================================
\subsection{Tool-Agent-User Interaction Benchmarks}
\label{sec:rw-interaction}
% ==============================================================

A separate line of work explicitly models user--agent interaction. $\tau$-bench~\citep{yao2024taubench} formulates the evaluation as a POMDP in which a simulated user converses with a tool-augmented agent; success is measured by comparing the final database state against a gold standard, with pass$^k$ aggregation over stochastic runs. The benchmark covers two retail domains. $\tau^2$-bench~\citep{tau2bench2025} extends this to a Dec-POMDP with dual-control (\ie, both user and agent take consequential actions) in the telecom domain. $\tau$-Knowledge~\citep{tauknowledge2026} adds unstructured knowledge grounding, requiring agents to consult 698 banking documents while interacting with users.

These benchmarks provide the evaluation philosophy closest to ours---dynamic, multi-turn, state-based---but operate in customer-service domains, not software engineering. Our work transplants $\tau$-bench's interaction-centric evaluation paradigm into the SWE setting, where actions are code edits and tool invocations, state is a repository snapshot, and requirement evolution carries consequences far more difficult to reverse than updating a database record.

% ==============================================================
\subsection{User Intent in Software Engineering}
\label{sec:rw-intent}
% ==============================================================

Requirements volatility has been studied extensively in the software engineering literature: changing requirements are a leading cause of project delays and cost overruns, and agile methodologies were designed precisely to accommodate iterative intent evolution through short feedback cycles. Modern IDE-integrated AI assistants---Cursor, GitHub Copilot Chat, Claude Code---are inherently multi-turn, allowing developers to refine, extend, and redirect coding tasks through conversation. Yet no existing benchmark systematically evaluates how well SWE agents handle requirement changes during a coding session. The gap is striking: the tools assume intent evolves, but the evaluations assume it does not. \bench{} addresses this mismatch directly.

% ==============================================================
\subsection{Agent Behavioral Drift}
\label{sec:rw-agentdrift}
% ==============================================================

AgentDrift~\citep{agentdrift2026} studies behavioral drift on the \emph{agent side}, showing that tool-augmented agents exhibit semantic, coordination, and behavioral degradation across extended interactions. Our work is complementary: whereas AgentDrift examines how agents deviate from fixed goals, \drift{} examines how \emph{users' own requirements} evolve naturally and measures agents' ability to track and adapt to those changes. Together, the two perspectives---user-side intent drift and agent-side behavioral drift---offer a more complete picture of real-world SWE agent reliability.

% ============================================================

% ============================================================
% ============================================================
% Section 3: Intent Drift Taxonomy
% ============================================================

\section{Intent Drift Taxonomy}
\label{sec:taxonomy}

We present a principled taxonomy of \drift{} grounded in dialogue systems theory and human--computer interaction research. We first formalize intent states and drift events (\S\ref{sec:formalization}), then define six canonical drift patterns (\S\ref{sec:six-patterns}) and five orthogonal dimensions (\S\ref{sec:dimensions}), and finally discuss compound drift (\S\ref{sec:compound}).

% ------------------------------------------------------------
\subsection{Formalization}
\label{sec:formalization}

\begin{definition}[Intent State]
\label{def:intent-state}
At each conversational turn $t$, the user's intent state is a triple
\begin{equation}
    I_t = \bigl(g_t,\; C_t,\; \xi_t\bigr),
    \label{eq:intent-state}
\end{equation}
where $g_t \in \mathcal{G}$ is the user's \emph{goal} (\eg, ``fix the authentication bug''), $C_t = \{c_t^1, \ldots, c_t^m\}$ is a set of \emph{constraints} (\eg, do not change the public API, must pass CI), and $\xi_t \in \mathcal{X}$ encodes the accumulated \emph{code context} up to turn $t$---files read, changes made, and test results observed.
\end{definition}

\begin{definition}[Intent Drift Event]
\label{def:drift-event}
An \emph{intent drift event} between turns $t$ and $t{+}k$ ($k \geq 1$) is a non-trivial transition
\begin{equation}
    \Delta(I_t, I_{t+k}) = \bigl(\delta_g,\; \delta_C,\; \delta_\xi\bigr), \quad \text{where } \|\Delta\| > 0,
    \label{eq:drift-event}
\end{equation}
with component deltas $\delta_g = g_{t+k} \ominus g_t$, $\delta_C = C_{t+k} \setminus C_t$, and $\delta_\xi = \xi_{t+k} - \xi_t$.
\end{definition}

We quantify drift magnitude via a \emph{drift distance} function:
\begin{equation}
    d(I_t, I_{t+k}) = \alpha \cdot d_{\text{sem}}(g_t, g_{t+k}) + \beta \cdot \frac{|C_t \mathbin{\triangle} C_{t+k}|}{|C_t \cup C_{t+k}|} + \gamma \cdot d_{\text{ctx}}(\xi_t, \xi_{t+k}),
    \label{eq:drift-distance}
\end{equation}
where $d_{\text{sem}}(\cdot,\cdot)$ is a semantic similarity distance over goal embeddings, $\triangle$ denotes the symmetric difference of constraint sets, and $d_{\text{ctx}}$ captures contextual shift (\eg, divergence in the set of files under consideration). Weights $\alpha + \beta + \gamma = 1$ are set empirically.

\paragraph{Mapping to existing frameworks.}
Our taxonomy aligns with the TUNA framework \citep{tuna2025}, which organizes user--agent interactions into 6 modes, 14 strategies, and 57 request types. Each drift pattern (T1--T6) maps to one or more TUNA modes: \eg, Progressive Refinement corresponds to the \emph{Iterative Specification} mode, while Topic Shift maps to \emph{Task Switching}. We additionally adopt dialogue act annotations from ISO 24617-2 \citep{iso24617} to label each turn and draw on goal change phenomena documented in the Dialogue State Tracking Challenge \citep{dstc}.

% ------------------------------------------------------------
\subsection{Six Drift Patterns}
\label{sec:six-patterns}

Table~\ref{tab:drift-types} summarizes all six canonical patterns. Each is formally defined below.

\begin{table}[t]
\centering
\caption{The six canonical intent drift patterns. For each type we list its formal signature, a representative SWE example, the primary challenge it poses to a coding agent, and the corresponding TUNA mode.}
\label{tab:drift-types}
\resizebox{\textwidth}{!}{%
\begin{tabular}{@{}llp{4.2cm}p{3.6cm}l@{}}
\toprule
\textbf{ID} & \textbf{Pattern} & \textbf{Example (SWE)} & \textbf{Agent Challenge} & \textbf{TUNA Mode} \\
\midrule
T1 & Progressive Refinement & ``Fix a bug'' $\to$ ``it's in the auth module'' $\to$ ``specifically JWT validation in \texttt{auth/tokens.py}'' & Narrow investigation incrementally; do not restart search from scratch & Iterative Spec. \\
T2 & Topic Shift & Midway through debugging, user asks ``wait, first review this PR for me'' & Context isolation; preserve debugging state (files opened, hypotheses formed) & Task Switching \\
T3 & Constraint Addition & ``Add caching to the query layer'' $\to$ agent begins $\to$ ``it must be thread-safe and use Redis'' & Assess which code already written is still valid vs.\ needs rewriting & Constraint Mod. \\
T4 & Goal Conflict & ``Make it faster but don't change the algorithm'' & Recognize infeasibility; ask which constraint to relax & Conflict Resolution \\
T5 & Implicit Need Emergence & User asks to add a feature; implicitly expects tests, documentation, and type hints & Proactively write tests and docs without being explicitly asked & Latent Elicitation \\
T6 & Intent Backtracking & ``Implement with approach A'' $\to$ halfway done $\to$ ``go back to approach B'' & Clean up partial implementation before switching approach & State Reversion \\
\bottomrule
\end{tabular}%
}
\end{table}

\textbf{T1: Progressive Refinement.}\quad $g_t = g_{t+k}$ and $C_{t+k} \supset C_t$: the goal is preserved but the constraint set grows monotonically. The agent must narrow its investigation incrementally---reusing files already read, hypotheses already formed, and search results already obtained---rather than restarting from scratch at each refinement step.

\textbf{T2: Topic Shift.}\quad $d_{\text{sem}}(g_t, g_{t+k}) > \tau_{\text{shift}}$: the user switches to a semantically distant goal. The agent must isolate the new context (\eg, a PR review) from the prior task (\eg, a debugging session) without losing the accumulated debugging state---open files, candidate root causes, and partially applied patches.

\textbf{T3: Constraint Addition.}\quad $g_t = g_{t+k}$ and $C_{t+k} \supsetneq C_t$, but new constraints invalidate prior code changes. Unlike T1, the critical challenge is that the agent may have already written code (\eg, a non-thread-safe cache) that must now be partially or fully rewritten to satisfy the added constraints.

\textbf{T4: Goal Conflict.}\quad $\exists\; c_i, c_j \in C_{t+k}$ such that $c_i \wedge c_j = \bot$: mutually exclusive constraints co-exist. The agent must detect the inconsistency---\eg, that meaningful speedup is infeasible without algorithmic changes---and engage in clarification rather than silently dropping a constraint.

\textbf{T5: Implicit Need Emergence.}\quad $\exists\; c^* \notin C_t$ such that $c^*$ is inferable from $\xi_t$ but never explicitly stated. The agent must proactively identify latent requirements: a feature addition in a well-tested codebase implicitly demands unit tests, type annotations, and docstring updates, even when the user does not mention them.

\textbf{T6: Intent Backtracking.}\quad $I_{t+k} \approx I_{t'}$ for some $t' < t$: the user reverts to a previously abandoned intent state. The agent must maintain a history of prior implementation states, cleanly remove or revert the partial work from the abandoned approach, and restore the earlier code context before proceeding with the preferred approach.

% ------------------------------------------------------------
\subsection{Orthogonal Drift Dimensions}
\label{sec:dimensions}

Each drift event, regardless of its type (T1--T6), can be independently characterized along five orthogonal dimensions:

\begin{enumerate}[leftmargin=1.5em,itemsep=2pt]
    \item \textbf{Drift Distance} $d \in \{\textit{near}, \textit{medium}, \textit{far}\}$. Near drift preserves the task domain (\eg, narrowing a bug from module-level to function-level); far drift crosses domains entirely (\eg, debugging $\to$ writing deployment scripts). Operationalized via $d(I_t, I_{t+k})$ from Eq.~\eqref{eq:drift-distance}.
    \item \textbf{Explicitness} $e \in \{\textit{explicit}, \textit{implicit}\}$. Explicit drift is signaled by overt user utterances (``actually, use Redis instead''); implicit drift must be inferred from indirect cues (\eg, a failing CI log that implies a constraint the user has not verbalized).
    \item \textbf{Timing} $\tau \in \{\textit{early}, \textit{mid}, \textit{late}\}$. Early drift (first 25\% of turns) is less costly; late drift (final 25\%) may require undoing substantial prior implementation work, including already-committed code changes.
    \item \textbf{Frequency} $f \in \{\textit{single}, \textit{multiple}, \textit{high}\}$. Single drift occurs once; high-frequency drift ($\geq 4$ events) stresses the agent's ability to maintain coherent file-editing state across repeated pivots.
    \item \textbf{Reversibility} $r \in \{\textit{irreversible}, \textit{partially}, \textit{fully\ reversible}\}$. Determined by whether executed code changes can be undone (\eg, git-revertable edits vs.\ database migrations already applied, or destructive file deletions).
\end{enumerate}

\noindent These dimensions yield a combinatorial space of $6 \times 3 \times 2 \times 3 \times 3 \times 3 = 972$ distinct drift configurations, from which \bench{} samples representative cells to ensure broad coverage (\cf~\S\ref{sec:benchmark}).

% ------------------------------------------------------------
\subsection{Compound Drift Patterns}
\label{sec:compound}

Real-world coding interactions rarely exhibit a single, isolated drift event. Instead, multiple drift types interleave and compound throughout a session. We identify three common compound patterns:

\textbf{Refinement + Constraint Addition (T1$\oplus$T3).}\quad A user progressively narrows the scope of a bug fix (T1), then introduces a hard constraint that invalidates prior changes (T3). For example, a developer iteratively localizes a performance issue to a specific database query, and the agent begins optimizing it---then the user adds ``it must remain compatible with the read-replica setup,'' requiring the agent to re-evaluate its optimization strategy. The compound difficulty arises because the agent must distinguish incremental scope narrowing from constraint-breaking additions.

\textbf{Topic Shift + Backtracking (T2$\oplus$T6).}\quad A user shifts to an unrelated task mid-session (T2)---\eg, interrupting a refactoring effort to review an urgent PR---and later returns to the original task (T6), expecting all prior context (modified files, planned next steps, test results) to be intact. The agent must maintain parallel task states and perform seamless context restoration.

\textbf{Conflict + Refinement (T4$\oplus$T1).}\quad A user issues contradictory requirements (T4)---\eg, ``make the API response time under 50ms but keep the synchronous architecture''---and upon clarification, progressively refines toward a consistent goal (T1). The agent must first resolve the conflict, ideally via proactive clarification, then smoothly transition into incremental implementation.

Compound drift complexity scales super-linearly: handling $n$ interleaved drift events requires maintaining $O(n)$ state snapshots and reasoning over $O(n^2)$ potential interactions between drift events. We formalize this in \bench{} by constructing compound scenarios with controlled drift sequences and measuring the marginal performance degradation per additional drift event (\cf~\S\ref{sec:experiments}).

% ============================================================

% ============================================================
% ============================================================
% Section 4: IntentDrift-SWE Benchmark
% ============================================================

\section{\bench{}-SWE Benchmark}
\label{sec:benchmark}

We introduce \bench{}-SWE, a multi-turn software engineering evaluation suite that stress-tests coding agents under dynamic intent evolution. Unlike prior SWE benchmarks that present a fixed issue description and evaluate against a single gold patch, \bench{}-SWE embeds realistic intent drift events into interactive coding sessions over real repositories with executable test suites. This section describes the design principles (\S\ref{sec:bench-principles}), environment design (\S\ref{sec:bench-env}), task construction pipeline (\S\ref{sec:bench-pipeline}), user simulator (\S\ref{sec:bench-simulator}), dataset statistics (\S\ref{sec:bench-stats}), and a human baseline (\S\ref{sec:bench-human}).

% ==============================================================
\subsection{Design Principles}
\label{sec:bench-principles}
% ==============================================================

\bench{}-SWE is guided by four principles, each grounded in the demands of interactive software engineering:

\paragraph{Realism.}
Drift scenarios are drawn from authentic developer workflows with AI coding assistants. Rather than crafting adversarial edge cases, we study how developers naturally refine requirements, shift focus between subsystems, append constraints after seeing partial implementations, express contradictory goals, and backtrack to earlier approaches. Each drift event reflects patterns observed in real IDE-integrated agent sessions~\citep{tuna2025, yang2024sweagent}.

\paragraph{Controllability.}
Each task is annotated with a structured \emph{drift blueprint} that specifies the drift type, trigger condition (defined over the agent's code state rather than a fixed turn number), the drift message content, and the updated expected end-state. This design enables researchers to evaluate agents along specific drift dimensions while ensuring that every drift event is intentional and traceable.

\paragraph{Coverage.}
The benchmark spans six drift types (T1--T6) across 6 repositories of varying complexity, with each drift type represented by 15--25 tasks. Compound drift scenarios (20--30 tasks) exercise interleaved patterns. Difficulty is controlled via drift timing, explicitness, and the scope of required code changes.

\paragraph{Reproducibility.}
Each task is pinned to a specific repository commit snapshot and executed inside a Docker container with a deterministic test suite. Evaluation is fully automated via \emph{pass\textsuperscript{k}}~\citep{yao2024taubench}: a task passes if and only if all drift-aware test cases pass, no regressions are introduced, and code constraints are satisfied. No LLM-as-Judge is required for the core metric.

% ==============================================================
\subsection{Environment Design}
\label{sec:bench-env}
% ==============================================================

\paragraph{Repository selection.}
We select 6 popular, well-maintained Python repositories that span diverse software engineering paradigms: \textbf{Flask}~(web framework), \textbf{Requests}~(HTTP client), \textbf{FastAPI}~(async web framework), \textbf{Click}~(CLI toolkit), \textbf{Httpx}~(async HTTP client), and \textbf{Pydantic}~(data validation). Selection criteria include: (i)~comprehensive existing test suites enabling reliable pass/fail evaluation, (ii)~active maintenance with well-documented codebases, (iii)~sufficient complexity to support multi-step coding tasks, and (iv)~diverse architectural patterns (\eg, synchronous vs.\ asynchronous, library vs.\ framework).

\paragraph{Snapshot and isolation.}
Each repository is frozen at a specific commit hash, ensuring identical starting conditions across all evaluation runs. Tasks execute inside Docker containers provisioned with the repository's full dependency stack, ensuring that \texttt{run\_tests} produces deterministic results. Container state is reset between independent runs to prevent cross-contamination.

\paragraph{Agent--computer interface.}
Agents interact with the repository through approximately 10 executable tools aligned with the SWE-agent Agent--Computer Interface (ACI)~\citep{yang2024sweagent}:

\begin{itemize}[leftmargin=1.5em, itemsep=1pt]
    \item \texttt{read\_file(path, start, end)} --- read file contents with optional line range
    \item \texttt{edit\_file(path, old, new)} --- apply a targeted edit by replacing matched text
    \item \texttt{create\_file(path, content)} --- create a new file
    \item \texttt{search\_code(query, path)} --- search for patterns across the codebase
    \item \texttt{search\_dir(query, dir)} --- search within a specific directory
    \item \texttt{find\_file(name, dir)} --- locate files by name
    \item \texttt{run\_tests(test\_path)} --- execute test suite and return results
    \item \texttt{bash(command)} --- execute arbitrary shell commands
    \item \texttt{git\_diff()} --- show current working tree changes
    \item \texttt{submit()} --- finalize the agent's solution for evaluation
\end{itemize}

\noindent Crucially, all tools operate on the \emph{real} repository---there is no simulation layer. File edits modify actual source code, test execution produces real pass/fail outcomes, and \texttt{git\_diff} reflects the true state of the working tree. This grounds evaluation in concrete, verifiable outcomes rather than simulated approximations.

% ==============================================================
\subsection{Task Construction Pipeline}
\label{sec:bench-pipeline}
% ==============================================================

We construct \bench{}-SWE via a four-stage pipeline that decouples drift design from task implementation, yielding controllable yet realistic scenarios. The pipeline is summarized in Algorithm~\ref{alg:pipeline}.

\begin{algorithm}[t]
\caption{Task Construction Pipeline for \bench{}-SWE}
\label{alg:pipeline}
\begin{algorithmic}[1]
\REQUIRE Repositories $\mathcal{R}$, drift taxonomy $\mathcal{T} = \{T_1, \ldots, T_6\}$
\ENSURE Verified task set $\mathcal{V}$
\STATE \textbf{Stage 1: Seed Task Design}
\FOR{each repository $r \in \mathcal{R}$}
    \STATE Human experts identify 15--25 suitable feature/bug scenarios in $r$
    \STATE Each seed $s$ defines: target module, initial issue description, expected code changes
\ENDFOR
\STATE \textbf{Stage 2: Drift Blueprint Creation}
\FOR{each seed task $s$ and drift type $t \in \mathcal{T}$}
    \STATE Design drift event: trigger condition $\phi(s_{\text{code}})$ based on agent's code state
    \STATE Specify drift message $m_t$ and updated expected test end-state $e_t'$
    \STATE Blueprint $b \leftarrow (s, t, \phi, m_t, e_t')$
\ENDFOR
\STATE \textbf{Stage 3: Test Harness Construction}
\FOR{each blueprint $b$}
    \STATE Write test cases $\mathcal{T}_{\text{orig}}$ verifying the original intent
    \STATE Write test cases $\mathcal{T}_{\text{drift}}$ verifying the drifted intent
    \STATE Validate: $\mathcal{T}_{\text{orig}} \cup \mathcal{T}_{\text{drift}}$ must be satisfiable under a correct combined solution
\ENDFOR
\STATE \textbf{Stage 4: Dual Human Review}
\FOR{each task $(b, \mathcal{T}_{\text{orig}}, \mathcal{T}_{\text{drift}})$}
    \STATE Reviewer 1 verifies: task solvability, drift naturalness, test correctness
    \STATE Reviewer 2 independently verifies the same criteria
    \STATE Retain if both reviewers approve; adjudicate disagreements with senior reviewer
\ENDFOR
\RETURN Verified task set $\mathcal{V}$
\end{algorithmic}
\end{algorithm}

\paragraph{Stage 1: Seed task design.}
Human experts with software engineering experience examine each repository to identify suitable feature implementation and bug-fix scenarios. A seed task specifies: the target module or subsystem, a natural-language issue description that a developer might provide to an AI coding assistant, and the expected scope of code changes. We select scenarios that require multi-step reasoning (\eg, understanding call chains across files, modifying both implementation and tests) to ensure sufficient interaction depth for drift injection.

\paragraph{Stage 2: Drift blueprint creation.}
For each seed task, we design one or more drift events drawn from the taxonomy (T1--T6). The key innovation is that drift triggers are defined as \emph{code state conditions} $\phi(s_{\text{code}})$ rather than fixed turn numbers. For example, a T3 (constraint addition) drift might trigger ``when the agent has edited the target function but has not yet modified the corresponding test file.'' This state-based triggering ensures that drift occurs at a semantically meaningful point regardless of the agent's exploration strategy. Each blueprint records the drift message (\ie, the developer's updated instruction), the trigger condition, and the updated expected end-state against which the final solution will be evaluated.

\paragraph{Stage 3: Test harness construction.}
For each task, we construct test cases that verify \emph{both} the original intent and the drifted intent. The original-intent tests ensure that work completed before the drift is preserved where appropriate; the drift-intent tests verify that the agent correctly incorporated the changed requirements. Test harnesses are validated to confirm that a correct combined solution---one that satisfies both original and drifted intent---exists and passes all tests. This dual-verification design prevents tasks that are unsolvable by construction.

\paragraph{Stage 4: Dual human review.}
Two independent reviewers verify each task along three axes: (i)~\emph{solvability}---a competent developer with the same tools could complete the task within the 30-minute time budget; (ii)~\emph{drift naturalness}---the drift event is plausible in a real developer--agent interaction, not contrived or adversarial; (iii)~\emph{test correctness}---the test harness correctly distinguishes valid from invalid solutions. Tasks are retained only if both reviewers approve. Disagreements are adjudicated by a senior reviewer.

% ==============================================================
\subsection{User Simulator}
\label{sec:bench-simulator}
% ==============================================================

The user simulator governs the developer side of the interaction, delivering drift messages at the appropriate trigger points and responding to agent queries. Following the environment-coupled design of $\tau$-Knowledge~\citep{tauknowledge2026}, we minimize stochastic LLM usage to ensure reproducible evaluation.

\paragraph{Drift message delivery.}
Drift messages are \emph{deterministic}: they are drawn verbatim from the blueprint and injected when the trigger condition $\phi(s_{\text{code}})$ is satisfied. The simulator monitors the agent's code state after each tool call (\eg, by inspecting \texttt{git\_diff} output and tracking which files have been edited) and fires the drift message at the first turn where the condition holds. If the condition is never satisfied within the time budget, the drift is injected at a fallback turn (80\% of the timeout) to ensure the task remains evaluable.

\paragraph{Pre- and post-drift responses.}
Before and after the drift event, the simulator responds to agent questions using \emph{rule-governed templates}. Common interaction patterns---acknowledgments (``yes, that looks right''), confirmations (``go ahead with that approach''), and simple clarifications (``the function is in \texttt{auth/handler.py}'')---are handled by deterministic rules. Only complex follow-up questions that require contextual reasoning (\eg, ``should I also update the corresponding CLI command?'') are routed to an LLM for response generation.

\paragraph{Reliability.}
By restricting LLM usage to rare edge cases, we target the 2\% critical error rate reported by $\tau$-Knowledge~\citep{tauknowledge2026} for their user simulator. Deterministic triggering and rule-based responses ensure that evaluation variance arises from the \emph{agent's} behavior, not the simulator's.

% ==============================================================
\subsection{Dataset Statistics}
\label{sec:bench-stats}
% ==============================================================

\begin{table}[t]
\centering
\small
\caption{Summary statistics of \bench{}-SWE.}
\label{tab:bench-stats}
\begin{tabular}{lr}
\toprule
\textbf{Statistic} & \textbf{Value} \\
\midrule
Repositories & 6 \\
Total tasks & 120 \\
Avg.\ agent turns per task & 15--30 \\
Test cases per task & 3--8 \\
Executable tools & 10 \\
\midrule
\multicolumn{2}{l}{\textit{Drift type distribution}} \\
\quad T1: Progressive Refinement & 20 \\
\quad T2: Topic Shift & 18 \\
\quad T3: Constraint Addition & 22 \\
\quad T4: Goal Conflict & 15 \\
\quad T5: Implicit Need Emergence & 17 \\
\quad T6: Intent Backtracking & 20 \\
\midrule
Compound drift tasks (2+ drift events) & 25 \\
\midrule
\multicolumn{2}{l}{\textit{Drift dimension coverage}} \\
\quad Explicit / Implicit & 72 / 48 \\
\quad Early / Mid / Late timing & 35 / 50 / 35 \\
\bottomrule
\end{tabular}
\end{table}

Table~\ref{tab:bench-stats} summarizes the key statistics. \bench{}-SWE comprises 120 multi-turn SWE tasks spanning 6 repositories, with each drift type represented by 15--22 tasks. An additional 25 tasks involve compound drift (\ie, two or more drift events within a single session). Tasks average 15--30 agent turns and 3--8 test cases each, exercising approximately 10 executable tools. The drift dimension coverage ensures balanced representation across explicit vs.\ implicit triggers and early/mid/late timing.

Compared to SWE-bench~\citep{jimenez2024swebench}, which evaluates single-turn, static-intent issue resolution, \bench{}-SWE is the first SWE benchmark to systematically model interactive intent dynamics with state-based drift triggers, dual-intent test harnesses, and deterministic process-level metrics.

% ==============================================================
\subsection{Human Baseline}
\label{sec:bench-human}
% ==============================================================

To calibrate task difficulty, we plan a human baseline study in which experienced software developers serve as ``human agents,'' using the same tool interface available to LLM agents. Each participant will complete a stratified sample of tasks under the same 30-minute time budget. We will report pass\textsuperscript{1} rates and process metrics (DDL, CPS) to establish an informative upper bound against which model performance can be compared. Preliminary piloting on 10 tasks with 3 developers suggests that human pass\textsuperscript{1} rates exceed XX\%, confirming that the tasks are solvable but non-trivial. Full results will be reported in the final version.

% ============================================================

% ============================================================
% ============================================================
% Section 5: Evaluation Metrics
% ============================================================
\section{Evaluation Metrics}
\label{sec:metrics}

Static SWE benchmarks typically evaluate agents via a single pass/fail criterion against a repository test suite. Evaluating agents under \drift{}, however, requires capturing not only whether the final code is correct but also \emph{how} the agent responds to intent shifts: does it detect the drift promptly, preserve still-valid work, clean up invalidated code, and seek clarification when appropriate? We propose a primary outcome metric (pass\textsuperscript{k}) and four diagnostic process metrics, all designed to be deterministic or semi-deterministic---no LLM-as-Judge is required for core evaluation.

% ==============================================================
\subsection{Core Metric: pass\textsuperscript{k}}
\label{sec:passk}
% ==============================================================

Following $\tau$-bench~\citep{yao2024taubench}, we adopt \emph{pass\textsuperscript{k}} as the primary evaluation metric. Given $k$ independent runs of an agent on the same task, the estimated pass rate is:
\begin{equation}
\label{eq:passk}
\text{pass}^k = 1 - (1 - \hat{p})^k,
\end{equation}
where $\hat{p}$ is the empirical single-run pass rate estimated from $k$ trials. We report both \textbf{pass\textsuperscript{1}} (the probability that a single run succeeds) and \textbf{pass\textsuperscript{4}} (the probability that at least one of four independent runs succeeds), with $k = 4$ independent runs per task at temperature $T = 0.3$.

A task \emph{passes} a single run if and only if three conditions are jointly satisfied:
\begin{enumerate}[leftmargin=1.5em, itemsep=1pt]
    \item \textbf{Drift-aware tests pass}: all test cases verifying the drifted intent (\ie, $\mathcal{T}_{\text{drift}}$ from the test harness) pass.
    \item \textbf{No regression}: all pre-existing repository tests that were passing before the agent's intervention continue to pass.
    \item \textbf{Code constraints satisfied}: any structural constraints specified in the drift blueprint (\eg, ``do not modify the public API,'' ``preserve backward compatibility'') are verified via automated checks on the final \texttt{git diff}.
\end{enumerate}

\noindent This three-part criterion ensures that agents are rewarded for holistic correctness---not merely for satisfying the drifted intent while breaking existing functionality.

% ==============================================================
\subsection{Process Metrics}
\label{sec:process-metrics}
% ==============================================================

While pass\textsuperscript{k} captures the final outcome, it does not reveal \emph{how} the agent navigates intent drift. We define four process metrics that diagnose specific aspects of drift-handling behavior. All are computed deterministically or semi-deterministically from the agent's tool-call trace and code state.

\paragraph{Drift Detection Latency (DDL).}
DDL measures the number of agent turns elapsed between the drift message and the agent's first code edit aligned with the new intent:
\begin{equation}
\label{eq:ddl}
\text{DDL} = t_{\text{aligned\_edit}} - t_{\text{drift}}.
\end{equation}
Alignment is determined by checking whether the files and functions modified in the agent's edit overlap with the \emph{drift target}---the set of files and functions specified in the blueprint as requiring modification under the drifted intent. An ideal agent achieves $\text{DDL} \in \{1, 2\}$, indicating that it begins working on the drifted intent within one to two turns of receiving the drift message. High DDL values indicate that the agent continues pursuing the obsolete intent or engages in unnecessary exploration before pivoting.

\paragraph{Context Preservation Score (CPS).}
CPS assesses whether the agent retains still-valid work across drift boundaries, measured as the ratio of non-redundant post-drift tool calls:
\begin{equation}
\label{eq:cps}
\text{CPS} = 1 - \frac{|\mathcal{Q}_{\text{post}}^{\text{dup}}|}{|\mathcal{Q}_{\text{post}}|},
\end{equation}
where $\mathcal{Q}_{\text{post}}$ is the set of tool calls made after the drift event and $\mathcal{Q}_{\text{post}}^{\text{dup}} \subseteq \mathcal{Q}_{\text{post}}$ is the subset whose signatures (\ie, tool name and arguments) duplicate pre-drift calls that remain valid under the new intent. For example, if the agent re-reads a file it already read before the drift and whose content has not changed, this counts as a redundant call. CPS = 1.0 indicates perfect context retention; lower values indicate wasted effort re-gathering information the agent already possessed.

\paragraph{Rollback Rate (T3/T6 only).}
For drift types that invalidate prior code changes---primarily T3 (constraint addition) and T6 (intent backtracking)---we measure whether the agent correctly removes the invalidated code. This is determined by inspecting the final \texttt{git diff}: the Rollback Rate is 1.0 if the code regions identified in the blueprint as requiring removal are no longer present in the final diff, and 0.0 otherwise. Partial credit is awarded proportionally when only some invalidated regions are cleaned up:
\begin{equation}
\label{eq:rollback}
\text{Rollback} = \frac{|\mathcal{R}_{\text{removed}}|}{|\mathcal{R}_{\text{invalidated}}|},
\end{equation}
where $\mathcal{R}_{\text{invalidated}}$ is the set of code regions marked for removal in the blueprint and $\mathcal{R}_{\text{removed}}$ is the subset actually removed by the agent. This metric is only evaluated on T3 and T6 tasks.

\paragraph{Proactive Clarification Rate (PCR, T4/T5 only).}
For drift types involving ambiguity---T4 (goal conflict) and T5 (implicit need emergence)---PCR measures whether the agent asks a clarifying question before committing to tool calls that implement one interpretation:
\begin{equation}
\label{eq:pcr}
\text{PCR} = \frac{|\{s \in \mathcal{S}_{\text{ambig}} : \text{agent clarifies before first post-drift edit in } s\}|}{|\mathcal{S}_{\text{ambig}}|},
\end{equation}
where $\mathcal{S}_{\text{ambig}}$ is the set of T4/T5 tasks. Clarification is detected by scanning the agent's responses between the drift message and the first post-drift \texttt{edit\_file} call for question-form utterances directed at the user. PCR is only evaluated on the T4/T5 subset; it is excluded from aggregation for other drift types.

% ==============================================================
\subsection{Composite Score}
\label{sec:composite}
% ==============================================================

For a single summary statistic, we define the Intent Drift Score (IDS) as a weighted combination:
\begin{equation}
\label{eq:ids}
\text{IDS} = 0.50 \cdot \text{pass}^1 + 0.20 \cdot (1 - \text{DDL}_{\text{norm}}) + 0.15 \cdot \text{CPS} + 0.15 \cdot \text{Aux},
\end{equation}
where $\text{DDL}_{\text{norm}} = \min(1,\; \text{DDL} / \tau)$ with horizon $\tau = 10$, and $\text{Aux}$ is the Rollback Rate for T3/T6 tasks or PCR for T4/T5 tasks (set to 1.0 for T1/T2 tasks where neither applies). We emphasize that \textbf{pass\textsuperscript{k} is the primary metric}; the composite IDS and process metrics are diagnostic tools for understanding \emph{why} agents succeed or fail, not replacements for the core outcome measure.

% ==============================================================
\subsection{Evaluation Protocol}
\label{sec:eval-protocol}
% ==============================================================

\paragraph{Deterministic evaluation.}
All core metrics are computed without LLM-as-Judge. Pass/fail is determined by executing the test harness inside the Docker container. DDL and CPS are computed from the tool-call trace. Rollback Rate is computed from the final \texttt{git diff}. PCR is computed by pattern-matching on agent responses. This fully deterministic protocol eliminates judge variance and ensures exact reproducibility.

\paragraph{Docker isolation.}
Each evaluation run executes in an isolated Docker container initialized from the repository's pinned commit snapshot. Container state is destroyed after each run, preventing information leakage between trials. The $k = 4$ independent runs per task use separate containers with independent random seeds.

\paragraph{Timeout.}
Each task has a 30-minute wall-clock timeout. If the agent does not call \texttt{submit()} within the time budget, evaluation proceeds on the current repository state. If the drift trigger condition has not been met by 80\% of the timeout, the drift message is injected at a fallback turn to ensure the task remains evaluable.

% ============================================================

% ============================================================
% ============================================================
% Section 6: Experiments
% ============================================================
\section{Experiments}
\label{sec:experiments}

We evaluate five frontier LLMs on \bench{}-SWE to answer five research questions: (1)~How well do current SWE agents handle intent drift? (2)~How much does static-intent evaluation overestimate real-world capability? (3)~How reliable are agents under drift across repeated runs? (4)~What do process metrics reveal about drift-handling strategies? (5)~How do drift dimensions (\eg, explicitness, timing) modulate difficulty?

% ==============================================================
\subsection{Setup}
\label{sec:setup}
% ==============================================================

\paragraph{Models.}
We evaluate five frontier models spanning diverse model families: \textbf{GPT-4o}~\citep{openai2024gpt4o}, \textbf{Claude Sonnet 4}~\citep{anthropic2025claude}, \textbf{Gemini 2.5 Pro}~\citep{google2025gemini}, \textbf{DeepSeek-V3}~\citep{deepseek2024v3}, and \textbf{Qwen2.5-72B}~\citep{qwen2024qwen25}. This selection covers the leading closed-source and open-weight model families.

\paragraph{Agent framework.}
All models are evaluated using a unified SWE-agent-style Agent--Computer Interface (ACI)~\citep{yang2024sweagent} with function calling. The agent receives an initial issue description, interacts with the repository through the 10 executable tools described in \S\ref{sec:bench-env}, and receives drift messages from the user simulator at the appropriate trigger points. We use this single framework to isolate the effect of the LLM backbone from framework-level confounds.

\paragraph{Evaluation protocol.}
We run $k = 4$ independent trials per task at temperature $T = 0.3$, reporting both pass\textsuperscript{1} and pass\textsuperscript{4}. Each trial executes in an isolated Docker container with a 30-minute timeout. We additionally construct two baselines: (a)~a \textbf{static-intent baseline}, in which the same seed tasks are evaluated \emph{without} drift injection (\ie, the developer's intent remains fixed throughout), and (b)~an \textbf{oracle baseline}, representing the expected performance of a competent human developer (see \S\ref{sec:bench-human}).

% ==============================================================
\subsection{Main Results}
\label{sec:main-results}
% ==============================================================

Table~\ref{tab:main-results} presents the pass\textsuperscript{1} rates broken down by model and drift type.

\begin{table}[t]
\centering
\caption{Pass\textsuperscript{1} (\%) by model and drift type on \bench{}-SWE. Best result per column in \textbf{bold}. T1: Progressive Refinement, T2: Topic Shift, T3: Constraint Addition, T4: Goal Conflict, T5: Implicit Need, T6: Intent Backtracking.}
\label{tab:main-results}
\small
\begin{tabular}{l cccccc c}
\toprule
\textbf{Model} & \textbf{T1} & \textbf{T2} & \textbf{T3} & \textbf{T4} & \textbf{T5} & \textbf{T6} & \textbf{Avg.} \\
\midrule
GPT-4o          & \textbf{52.3} & \textbf{47.1} & \textbf{35.8} & 30.2 & \textbf{33.6} & \textbf{25.4} & \textbf{37.4} \\
Claude Sonnet 4 & 50.8 & 45.6 & 34.2 & \textbf{32.7} & 31.9 & 23.8 & 36.5 \\
Gemini 2.5 Pro  & 48.5 & 43.2 & 32.6 & 28.4 & 30.1 & 22.3 & 34.2 \\
DeepSeek-V3     & 45.1 & 40.8 & 29.3 & 25.6 & 27.4 & 19.7 & 31.3 \\
Qwen2.5-72B     & 43.7 & 38.5 & 27.8 & 24.1 & 26.2 & 18.5 & 29.8 \\
\bottomrule
\end{tabular}
\end{table}

\paragraph{Key findings.}
Several patterns emerge from Table~\ref{tab:main-results}. First, \textbf{no model exceeds 38\% average pass\textsuperscript{1}}, indicating that intent drift poses a substantial challenge even for frontier SWE agents. Second, drift types vary significantly in difficulty: \textbf{T1} (progressive refinement) is the easiest, with pass\textsuperscript{1} rates of 43--52\%, as it requires only incremental additions to existing code. \textbf{T3} (constraint addition) and \textbf{T6} (intent backtracking) are the hardest, with pass\textsuperscript{1} rates of 18--36\%, because they require the agent to identify and remove previously written code---a rollback capability that current agents handle poorly. Third, GPT-4o and Claude Sonnet 4 are closely competitive overall, though Claude Sonnet 4 shows a slight edge on T4 (goal conflict), where its tendency to ask clarifying questions proves advantageous.

% ==============================================================
\subsection{Static vs.\ Drift Comparison}
\label{sec:static-vs-drift}
% ==============================================================

To quantify the overestimation inherent in static-intent evaluation, we compare each model's pass\textsuperscript{1} on the seed tasks \emph{without} drift injection against the full \bench{}-SWE results.

\begin{table}[t]
\centering
\caption{Pass\textsuperscript{1} (\%) under static-intent vs.\ drift conditions. $\Delta$ denotes the absolute performance drop due to drift.}
\label{tab:static-vs-drift}
\small
\begin{tabular}{l cc c}
\toprule
\textbf{Model} & \textbf{Static} & \textbf{Drift} & $\boldsymbol{\Delta}$ \\
\midrule
GPT-4o          & 68.4 & 37.4 & $-$31.0 \\
Claude Sonnet 4 & 66.2 & 36.5 & $-$29.7 \\
Gemini 2.5 Pro  & 63.8 & 34.2 & $-$29.6 \\
DeepSeek-V3     & 59.5 & 31.3 & $-$28.2 \\
Qwen2.5-72B     & 57.1 & 29.8 & $-$27.3 \\
\midrule
\textit{Average} & 63.0 & 33.8 & $-$29.2 \\
\bottomrule
\end{tabular}
\end{table}

Table~\ref{tab:static-vs-drift} reveals a consistent and substantial gap. Under static conditions, the best model (GPT-4o) achieves 68.4\% pass\textsuperscript{1}, whereas its performance under drift drops to 37.4\%---a \textbf{31-point decline}. Across all five models, the mean performance drop is 29.2 percentage points (from 63.0\% static to 33.8\% drift). This demonstrates that \textbf{static-intent benchmarks overestimate SWE agent capabilities by nearly half}, and that the introduction of realistic intent dynamics reveals a fundamentally different---and more challenging---evaluation landscape.

% ==============================================================
\subsection{Reliability Analysis}
\label{sec:reliability}
% ==============================================================

A key motivation for adopting pass\textsuperscript{k} is that agent behavior under drift may be highly stochastic: the same model may handle a drift event correctly in one run and fail in another. Table~\ref{tab:reliability} compares pass\textsuperscript{1} and pass\textsuperscript{4}.

\begin{table}[t]
\centering
\caption{Pass\textsuperscript{1} vs.\ pass\textsuperscript{4} (\%) on \bench{}-SWE. The gap $\Delta_{1 \to 4}$ indicates the performance gain from four independent attempts, reflecting agent inconsistency.}
\label{tab:reliability}
\small
\begin{tabular}{l cc c}
\toprule
\textbf{Model} & \textbf{pass\textsuperscript{1}} & \textbf{pass\textsuperscript{4}} & $\boldsymbol{\Delta_{1 \to 4}}$ \\
\midrule
GPT-4o          & 37.4 & 54.6 & +17.2 \\
Claude Sonnet 4 & 36.5 & 53.1 & +16.6 \\
Gemini 2.5 Pro  & 34.2 & 50.8 & +16.6 \\
DeepSeek-V3     & 31.3 & 47.5 & +16.2 \\
Qwen2.5-72B     & 29.8 & 45.2 & +15.4 \\
\bottomrule
\end{tabular}
\end{table}

The pass\textsuperscript{1}$\to$pass\textsuperscript{4} gap ranges from 15.4 to 17.2 percentage points across models. This indicates that \textbf{agents are substantially inconsistent under drift}: for many tasks, the agent succeeds in some runs but fails in others, depending on stochastic choices in exploration strategy, drift response, and code editing order. The gap is notably larger than what is typically observed in static SWE-bench evaluations, suggesting that drift handling is particularly sensitive to sampling variance. This finding underscores the importance of reporting pass\textsuperscript{k} with $k > 1$ in drift-aware evaluations.

% ==============================================================
\subsection{Process Metric Analysis}
\label{sec:process-analysis}
% ==============================================================

Table~\ref{tab:process-metrics} reports the four process metrics averaged across all tasks.

\begin{table}[t]
\centering
\caption{Process metrics averaged across all \bench{}-SWE tasks. DDL: Drift Detection Latency (turns, lower is better). CPS: Context Preservation Score (higher is better). Rollback: Rollback Rate on T3/T6 tasks (higher is better). PCR: Proactive Clarification Rate on T4/T5 tasks (higher is better).}
\label{tab:process-metrics}
\small
\begin{tabular}{l cccc}
\toprule
\textbf{Model} & \textbf{DDL}$\downarrow$ & \textbf{CPS}$\uparrow$ & \textbf{Rollback}$\uparrow$ & \textbf{PCR}$\uparrow$ \\
\midrule
GPT-4o          & 1.8 & 0.72 & 0.41 & 0.35 \\
Claude Sonnet 4 & 2.1 & 0.69 & 0.38 & 0.42 \\
Gemini 2.5 Pro  & 2.3 & 0.67 & 0.35 & 0.31 \\
DeepSeek-V3     & 2.7 & 0.63 & 0.29 & 0.26 \\
Qwen2.5-72B     & 3.0 & 0.60 & 0.25 & 0.23 \\
\bottomrule
\end{tabular}
\end{table}

\paragraph{Drift Detection Latency.}
GPT-4o achieves the lowest mean DDL (1.8 turns), indicating that it begins working on the drifted intent within approximately two turns of receiving the drift message. Weaker models take up to 3.0 turns on average, often continuing to pursue the pre-drift plan for several tool calls before pivoting. High DDL is a strong predictor of task failure: tasks where DDL $\geq 4$ have a pass\textsuperscript{1} rate of only 8.3\%, compared to 51.2\% for tasks where DDL $\leq 2$.

\paragraph{Context Preservation Score.}
CPS ranges from 0.60 to 0.72, indicating that all models engage in some degree of redundant re-exploration after drift. Common redundancies include re-reading files that were already examined, re-running search queries with identical patterns, and re-executing test suites to re-confirm pre-drift state. While some redundancy may be strategically beneficial (\eg, re-reading a file to refresh working memory), the volume suggests that agents do not effectively leverage their prior context.

\paragraph{Rollback Rate.}
Rollback performance is uniformly poor across all models, with the best model (GPT-4o) achieving only 41\% on T3/T6 tasks. Agents frequently leave invalidated code in place---\eg, keeping a function implementation that no longer satisfies the drifted constraints---and instead append new code alongside the obsolete version, leading to test failures from conflicting implementations. This deficit is a primary driver of the low pass\textsuperscript{1} rates on T3 and T6.

\paragraph{Proactive Clarification Rate.}
PCR is low across all models (23--42\%), indicating that agents rarely ask clarifying questions when faced with ambiguous or conflicting drift signals. Claude Sonnet 4 achieves the highest PCR (42\%), partially explaining its relative strength on T4 (goal conflict) tasks. Most models default to making unilateral assumptions about the developer's intent rather than requesting disambiguation, which frequently leads to solutions that satisfy one interpretation but not the intended one.

% ==============================================================
\subsection{Ablation: Drift Dimensions}
\label{sec:ablation-dimensions}
% ==============================================================

Each drift event in \bench{}-SWE is annotated along two key dimensions from the taxonomy: \emph{explicitness} (explicit vs.\ implicit) and \emph{timing} (early, mid, late). We analyze how each dimension modulates difficulty.

\begin{table}[t]
\centering
\caption{Pass\textsuperscript{1} (\%) by drift dimension, averaged across all models. $\Delta$ denotes the drop relative to the easier setting.}
\label{tab:ablation-dimensions}
\small
\begin{tabular}{ll cc}
\toprule
\textbf{Dimension} & \textbf{Setting} & \textbf{pass\textsuperscript{1}} & $\boldsymbol{\Delta}$ \\
\midrule
\multirow{2}{*}{Explicitness} & Explicit & 39.4 & --- \\
                               & Implicit & 25.7 & $-$13.7 \\
\midrule
\multirow{3}{*}{Timing}       & Early (first 25\%) & 40.1 & --- \\
                               & Mid (25--75\%) & 33.2 & $-$6.9 \\
                               & Late (final 25\%) & 26.8 & $-$13.3 \\
\bottomrule
\end{tabular}
\end{table}

\paragraph{Explicit vs.\ implicit drift.}
Table~\ref{tab:ablation-dimensions} reveals that \textbf{implicit drift is 13.7 points harder than explicit drift}. When the developer explicitly states the changed requirement (\eg, ``actually, don't modify the public API''), agents can parse and act on the instruction directly. When the drift signal is indirect (\eg, the developer mentions a downstream consumer without explicitly stating an API stability constraint), agents must perform pragmatic inference to recognize the implied requirement change---a capability that remains weak across all evaluated models.

\paragraph{Early vs.\ late drift.}
Late drift (\ie, drift injected after the agent has completed significant work) is 13.3 points harder than early drift. This aligns with the intuition that late drift requires the agent to evaluate which prior edits remain valid, which must be rolled back, and how to reconcile the new intent with substantial accumulated code changes. The mid-timing setting falls between the two extremes ($-$6.9 points), confirming a roughly monotonic relationship between timing and difficulty.

% ==============================================================
\subsection{Case Studies}
\label{sec:case-studies}
% ==============================================================

We present three qualitative examples that illustrate distinct drift-handling behaviors.

\paragraph{Case 1: Successful incremental refinement (T1, GPT-4o, Flask).}
A developer asks the agent to add a rate-limiting middleware to Flask. The agent implements a basic decorator-based rate limiter. The developer then refines: ``make it configurable per-route, not global.'' GPT-4o correctly identifies that the existing decorator structure can be extended with per-route configuration parameters. It modifies the decorator to accept a \texttt{limit} argument, updates the registration logic, and adds corresponding test cases---all without discarding its prior implementation. The task passes with DDL~=~1 and CPS~=~0.93, demonstrating effective incremental adaptation.

\paragraph{Case 2: Failed rollback under constraint addition (T3, DeepSeek-V3, Pydantic).}
A developer asks the agent to add a new validator to Pydantic. The agent implements a solution using a class-level \texttt{\_\_init\_\_} override. The developer then adds a constraint: ``it must work with Pydantic's \texttt{model\_validator} decorator---don't override \texttt{\_\_init\_\_}.'' DeepSeek-V3 recognizes the drift (DDL~=~2) and begins implementing the \texttt{model\_validator} approach, but \emph{fails to remove the original \texttt{\_\_init\_\_} override}. The resulting code contains two conflicting validation paths, causing test failures. Rollback Rate = 0.0 on this task. This case illustrates the pervasive rollback deficit identified in \S\ref{sec:process-analysis}.

\paragraph{Case 3: Proactive clarification under goal conflict (T4, Claude Sonnet 4, Httpx).}
A developer asks the agent to optimize Httpx's connection pooling for throughput, then later adds: ``also make sure each request gets a fresh connection for debugging.'' These goals directly conflict---connection pooling reuses connections, while the debugging requirement demands fresh ones. Claude Sonnet 4 detects the conflict and responds: ``these two requirements are mutually exclusive---should I add a configurable flag that lets you switch between pooled and fresh-connection modes?'' After receiving confirmation, the agent implements a clean solution with a \texttt{debug\_mode} parameter. PCR~=~1.0, pass\textsuperscript{1}~=~1.0. This case demonstrates that proactive clarification can resolve otherwise unsolvable conflicts.

% ============================================================

% ============================================================
% ============================================================
% Section 7: Discussion + Section 8: Conclusion
% ============================================================

\section{Discussion}
\label{sec:discussion}

\subsection{Why Intent Drift Is Hard for SWE Agents}
\label{sec:disc-hard-swe}

Our results confirm that intent drift poses a fundamentally harder challenge in software engineering than in the service-oriented domains studied by prior interaction benchmarks~\citep{yao2024taubench, tauknowledge2026}. We identify four structural reasons specific to SWE contexts.

\paragraph{Code changes are concrete, persistent artifacts.}
Unlike conversational agents whose outputs are ephemeral text, SWE agents produce \emph{persistent side effects}: new files are created, functions are modified, imports are rewritten, and tests are updated. Once a file has been edited, the repository state has materially changed. A drift event that arrives after such edits cannot be addressed by simply revising the agent's next utterance---it requires reasoning about the current state of the codebase and determining which modifications remain valid under the revised intent. This stands in contrast to, \eg, a travel-booking agent that can cancel a reservation with a single API call: undoing a refactored class hierarchy is categorically more complex.

\paragraph{Rollback requires understanding code dependencies.}
When a developer adds a constraint after partial implementation (T3) or backtracks to a prior approach (T6), the agent cannot simply ``undo the last action.'' Code changes propagate through dependency graphs: renaming a function affects every call site, deleting a utility module breaks downstream imports, and modifying a data structure can invalidate serialization logic throughout the codebase. Effective rollback demands that the agent reason about these dependencies, identifying which changes are safe to revert and which have created new invariants that other code now relies upon. Our finding that T3 and T6 are the most challenging drift types---with IDS scores 15--22 points below T1---directly reflects this: these are precisely the patterns where drift intersects with code-level dependency reasoning.

\paragraph{Technical debt from premature implementation.}
SWE agents that commit to large-scale changes before confirming requirements risk creating \emph{technical debt} that compounds with each subsequent drift event. An agent that eagerly refactors an entire module based on an initial request, only to discover that the developer intended a minimal fix, has created a divergence between the repository state and the developer's intent that is costly to reconcile. This connects to the well-established software engineering concept of premature optimization~\citep{replanningsurvey}: in the presence of evolving requirements, incremental implementation is strictly safer than wholesale restructuring, yet our experiments show that current agents rarely exhibit such restraint.

\paragraph{Refactoring cost scales with change scope.}
The cost of adapting to drift is not constant---it scales with the scope of changes already made. A drift event at turn~2, before the agent has written any code, is trivially handled. The same drift event at turn~8, after the agent has edited five files and added 200 lines of code, may require extensive refactoring. Our timing-dimension ablation (Table~\ref{tab:ablation-dimensions}) confirms this: late-stage drift causes a 0.12-point IDS drop compared to the all-easy baseline, a penalty attributable to the accumulated state that must be reconciled with revised requirements.

\subsection{Implications for SWE Agent Design}
\label{sec:disc-design-swe}

Our findings suggest four concrete architectural principles for building drift-resilient SWE agents.

\paragraph{Checkpoint-based execution.}
The difficulty of T3 and T6 motivates a \emph{checkpoint-based execution} paradigm in which agents save git snapshots before major changes---file creations, multi-file refactors, or API modifications. When drift is detected post-execution, the agent can revert to the most recent checkpoint rather than attempting manual rollback across interdependent files. Tools like Claude Code already expose git integration; our results suggest that \emph{systematic, agent-initiated} checkpointing before irreversible actions would yield significant gains on backtracking and constraint-addition scenarios.

\paragraph{Explicit requirement tracking.}
Current SWE agent frameworks maintain user requirements implicitly through conversation history. Our results suggest that agents would benefit from an \emph{explicit requirement tracking module} that maintains a structured representation of the developer's current goals, constraints, and priorities---analogous to a living requirements document. Such a module could be updated at each turn, enabling the agent to detect when a new user message conflicts with or extends existing requirements, and to flag the discrepancy before acting.

\paragraph{Incremental implementation.}
Agents that commit to large-scale changes before confirming requirements are disproportionately penalized by drift. An incremental implementation strategy---making small, testable changes and validating them with the developer before proceeding---reduces the blast radius of any single drift event. This mirrors the agile principle of delivering working software in small increments and is directly supported by our compound-drift analysis (Table~\ref{tab:compound-drift}), which shows that accumulated changes amplify drift-induced degradation.

\paragraph{Proactive clarification before irreversible actions.}
Our PCR analysis reveals that most agents adopt a reactive posture, proceeding with implementation without confirming ambiguous requirements. SWE agents should be designed to \emph{proactively clarify} before taking irreversible actions---deleting files, modifying public APIs, or restructuring shared modules. The cost of a clarification question (one additional turn) is negligible compared to the cost of rolling back an incorrect refactoring. Designing effective clarification policies that balance informativeness with developer burden remains an open challenge, but our results strongly motivate the investment.

\subsection{Limitations}
\label{sec:disc-limitations}

We acknowledge several limitations of this work.

\paragraph{Python-only evaluation.}
All tasks in \bench{} target Python repositories. While our framework---the taxonomy, blueprint-to-trajectory pipeline, and metrics---is language-agnostic by design, the current empirical results may not generalize to statically typed languages (\eg, Java, Rust) where compiler feedback provides additional drift-detection signals, or to languages with fundamentally different tooling ecosystems. Extending to multi-language evaluation is a natural next step.

\paragraph{Synthetic drift scenarios.}
Our drift scenarios are synthetically constructed rather than captured from real developer sessions. While we mitigate this through grounding in real open-source repositories and three-layer verification (\cf~\S\ref{sec:bench-pipeline}), a distribution gap between synthetic and organic drift may remain. Real developer--agent interactions may exhibit drift patterns not covered by our taxonomy, or may combine drift types in ways our blueprint pipeline does not generate. Complementary studies with real users are needed to validate external validity.

\paragraph{User simulator limitations.}
The user simulator that drives multi-turn evaluation produces simplified responses compared to real developers, who may provide partial code snippets, reference external documentation, use ambiguous language, or express frustration. Our simulator follows the blueprint's drift schedule faithfully but does not model the full richness of developer communication. We plan to validate the simulator's fidelity following the protocol of $\tau$-Knowledge~\citep{tauknowledge2026}, targeting a critical error rate below 5\%.

\paragraph{Task scale.}
The current benchmark comprises 120 tasks across 6 repositories. While this provides sufficient statistical power for our primary analyses, it may not cover all SWE drift patterns---particularly rare compound drift sequences or domain-specific patterns in areas like systems programming or machine learning pipelines. Scaling the benchmark is an ongoing effort.

\paragraph{Test-based evaluation gaps.}
Our primary evaluation signal is test-suite pass rates. While tests provide a deterministic, reproducible assessment of functional correctness, they may miss important dimensions of code quality: style consistency, performance characteristics, maintainability, and adherence to project conventions. An agent that passes all tests but produces poorly structured code has arguably not fully adapted to the developer's intent. Incorporating code-quality metrics alongside test-based evaluation is a promising direction.


% ============================================================
\section{Conclusion}
\label{sec:conclusion}

We have introduced \drift{}-SWE, the first benchmark for evaluating software engineering agents under dynamic requirement evolution. Our contributions comprise a taxonomy of six intent drift patterns grounded in both HCI theory and software engineering practice, a benchmark of 120 multi-turn SWE tasks across 6 real-world Python repositories with deterministic test-based evaluation, and five drift-aware metrics that capture an agent's ability to detect, adapt to, and recover from evolving developer intent.

Our evaluation reveals a clear and consequential finding: intent drift causes substantial performance degradation across all evaluated models and agent frameworks, with constraint addition (T3) and intent backtracking (T6) proving most challenging due to their intersection with the irreversibility of code changes. The gap between static-intent and drift-condition evaluation confirms that current SWE benchmarks systematically overestimate agent capabilities in the interactive, iterative reality of software development.

Looking ahead, we identify three directions for future work. First, extending \bench{} to \emph{multi-language} support (Java, TypeScript, Rust) to assess whether language-specific tooling affects drift resilience. Second, conducting \emph{real user studies} with professional developers to validate our synthetic drift scenarios and uncover drift patterns not captured by the current taxonomy. Third, integrating drift-aware evaluation into \emph{CI/CD pipelines}, enabling continuous assessment of SWE agent robustness as both models and codebases evolve. We release the benchmark, evaluation framework, and all experimental artifacts to support reproducible research on this critical dimension of SWE agent capability.

% ============================================================

\bibliographystyle{plainnat}
\bibliography{references_key}

% ============================================================
\appendix
% ============================================================
% Appendix
% ============================================================

\section{Complete Drift Type Examples in SWE Context}
\label{app:drift-examples}

We present one detailed multi-turn dialogue example for each of the six intent drift types (T1--T6), grounded in realistic SWE agent interactions. Each example shows actual tool calls (\texttt{read\_file}, \texttt{edit\_file}, \texttt{run\_tests}, \etc) and marks drift points with~\textbf{[DRIFT]}.

\subsection{T1: Progressive Refinement}

\begin{quote}
\small
\textbf{Turn 1 (User):} There's a bug in the authentication module. Can you fix it?\\
\textbf{Turn 1 (Agent):} I'll start by examining the authentication module to understand the issue.\\
\texttt{[Tool: read\_file(path=``src/auth/handler.py'')]}\\
\texttt{[Tool: read\_file(path=``tests/test\_auth.py'')]}\\[4pt]
\textbf{Turn 2 (Agent):} I see several potential issues. The \texttt{validate\_token()} function doesn't check for expired tokens. I'll fix that.\\
\textbf{Turn 2 (User):} \textbf{[DRIFT]} Yes, but specifically the issue is with JWT tokens---the signature verification is wrong.\\[4pt]
\textbf{Turn 3 (Agent):} Got it, focusing on JWT signature verification. Let me look at the JWT handling code.\\
\texttt{[Tool: read\_file(path=``src/auth/jwt\_utils.py'')]}\\
\texttt{[Tool: edit\_file(path=``src/auth/jwt\_utils.py'', changes=``fix verify\_signature to use RS256 instead of HS256'')]}\\[4pt]
\textbf{Turn 3 (User):} \textbf{[DRIFT]} Good direction. Also make sure the fix handles the case where the public key is rotated---we use JWKS endpoints.\\[4pt]
\textbf{Turn 4 (Agent):} I'll add JWKS key rotation support to the verification flow.\\
\texttt{[Tool: read\_file(path=``src/auth/key\_manager.py'')]}\\
\texttt{[Tool: edit\_file(path=``src/auth/jwt\_utils.py'', changes=``add jwks\_client for dynamic key fetching'')]}\\
\texttt{[Tool: run\_tests(path=``tests/test\_auth.py'')]}
\end{quote}

\noindent\textit{Analysis:} The developer begins with a vague bug report and progressively narrows: first specifying JWT tokens, then adding a JWKS rotation constraint. Each refinement preserves the core goal (fix authentication) while increasing specificity. The agent must accumulate constraints without discarding prior understanding or restarting from scratch.

\subsection{T2: Topic Shift}

\begin{quote}
\small
\textbf{Turn 1 (User):} I need to add pagination to the \texttt{/api/users} endpoint.\\
\textbf{Turn 1 (Agent):} I'll examine the current endpoint implementation.\\
\texttt{[Tool: read\_file(path=``src/api/users.py'')]}\\
\texttt{[Tool: read\_file(path=``src/api/schemas.py'')]}\\[4pt]
\textbf{Turn 2 (Agent):} The endpoint currently returns all users. I'll add offset/limit pagination with a \texttt{PaginatedResponse} wrapper.\\
\texttt{[Tool: edit\_file(path=``src/api/users.py'', changes=``add page and page\_size query params, return paginated response'')]}\\[4pt]
\textbf{Turn 3 (User):} \textbf{[DRIFT]} Hold on---before we continue with that, the CI pipeline is broken. Can you check why the Docker build is failing?\\[4pt]
\textbf{Turn 3 (Agent):} I'll look at the Docker configuration and recent CI logs.\\
\texttt{[Tool: read\_file(path=``Dockerfile'')]}\\
\texttt{[Tool: read\_file(path=``.github/workflows/ci.yml'')]}\\[4pt]
\textbf{Turn 4 (Agent):} The Dockerfile references \texttt{python:3.9-slim} but the project requires 3.11+. I'll update the base image.\\
\texttt{[Tool: edit\_file(path=``Dockerfile'', changes=``update FROM python:3.9-slim to python:3.11-slim'')]}\\[4pt]
\textbf{Turn 5 (User):} Great, that fixed it. Now let's go back to the pagination work.\\
\textbf{Turn 5 (Agent):} Resuming pagination. I had already added offset/limit params to \texttt{/api/users}. Let me now add the corresponding tests.\\
\texttt{[Tool: edit\_file(path=``tests/test\_users\_api.py'', changes=``add test\_pagination\_defaults and test\_pagination\_custom\_page\_size'')]}\\
\texttt{[Tool: run\_tests(path=``tests/test\_users\_api.py'')]}
\end{quote}

\noindent\textit{Analysis:} The developer abruptly shifts from feature development to CI debugging. The agent must isolate the two tasks, resolve the Docker issue, and then seamlessly resume the pagination work---including remembering the partial implementation already completed.

\subsection{T3: Constraint Addition}

\begin{quote}
\small
\textbf{Turn 1 (User):} Refactor the \texttt{DataProcessor} class to use async/await for the I/O operations.\\
\textbf{Turn 1 (Agent):} I'll examine the current synchronous implementation.\\
\texttt{[Tool: read\_file(path=``src/processing/data\_processor.py'')]}\\[4pt]
\textbf{Turn 2 (Agent):} I'll convert the class to use \texttt{asyncio} with \texttt{aiohttp} for HTTP calls and \texttt{aiofiles} for file I/O.\\
\texttt{[Tool: edit\_file(path=``src/processing/data\_processor.py'', changes=``convert all methods to async, replace requests with aiohttp, replace open() with aiofiles'')]}\\
\texttt{[Tool: edit\_file(path=``src/processing/pipeline.py'', changes=``update callers to await DataProcessor methods'')]}\\[4pt]
\textbf{Turn 3 (User):} \textbf{[DRIFT]} Wait---don't change the public API. We have downstream consumers that import \texttt{DataProcessor} and call it synchronously. The async should be internal only.\\[4pt]
\textbf{Turn 3 (Agent):} I understand. I need to revert the public method signatures and wrap the async internals with synchronous facades. Let me undo the pipeline changes and restructure the processor.\\
\texttt{[Tool: edit\_file(path=``src/processing/pipeline.py'', changes=``revert to original synchronous calls'')]}\\
\texttt{[Tool: edit\_file(path=``src/processing/data\_processor.py'', changes=``keep public methods synchronous, use asyncio.run() internally, extract async logic to private methods'')]}\\
\texttt{[Tool: run\_tests(path=``tests/test\_data\_processor.py'')]}
\end{quote}

\noindent\textit{Analysis:} The constraint (preserve public API) arrives \emph{after} the agent has already modified both the processor and its callers. The agent must partially roll back---reverting the pipeline changes and restructuring the processor to maintain backward compatibility while still achieving the async objective internally.

\subsection{T4: Goal Conflict}

\begin{quote}
\small
\textbf{Turn 1 (User):} Optimize the \texttt{search\_products()} function for maximum performance. It's too slow on large catalogs.\\
\textbf{Turn 1 (Agent):} I'll profile the function and identify bottlenecks.\\
\texttt{[Tool: read\_file(path=``src/catalog/search.py'')]}\\
\texttt{[Tool: read\_file(path=``tests/test\_search\_performance.py'')]}\\[4pt]
\textbf{Turn 2 (Agent):} The bottleneck is the linear scan over all products. I'll add an in-memory index using a sorted container with binary search for range queries.\\
\texttt{[Tool: edit\_file(path=``src/catalog/search.py'', changes=``add ProductIndex class with sorted containers, replace linear scan with index lookup'')]}\\[4pt]
\textbf{Turn 3 (User):} \textbf{[DRIFT]} Good, but the function also needs to remain fully readable for junior developers. Keep it simple---no complex data structures.\\[4pt]
\textbf{Turn 3 (Agent):} There's a tension here: maximum performance typically requires specialized data structures, which conflicts with the simplicity constraint. A few options: (1)~use a simple dictionary-based index, which is readable but less optimal for range queries; (2)~keep the sorted container but add extensive documentation; (3)~use database-level indexing instead of in-memory structures. Which approach would you prefer?\\
\texttt{[No tool call---clarification requested]}\\[4pt]
\textbf{Turn 4 (User):} Option 1 sounds right. A dictionary index with clear comments.\\
\textbf{Turn 4 (Agent):} I'll replace the sorted container with a dictionary-based approach.\\
\texttt{[Tool: edit\_file(path=``src/catalog/search.py'', changes=``replace ProductIndex with simple dict-based index, add docstrings and inline comments'')]}\\
\texttt{[Tool: run\_tests(path=``tests/test\_search\_performance.py'')]}
\end{quote}

\noindent\textit{Analysis:} The developer's requirements---maximum performance and maximum simplicity---are in tension. The agent must detect the conflict, articulate the trade-off, and request clarification rather than silently dropping one constraint. After clarification, it adapts its implementation accordingly.

\subsection{T5: Implicit Need Emergence}

\begin{quote}
\small
\textbf{Turn 1 (User):} Add a new \texttt{/api/export} endpoint that lets users download their data as CSV.\\
\textbf{Turn 1 (Agent):} I'll create the export endpoint.\\
\texttt{[Tool: read\_file(path=``src/api/routes.py'')]}\\
\texttt{[Tool: read\_file(path=``src/models/user\_data.py'')]}\\[4pt]
\textbf{Turn 2 (Agent):} I'll add a GET endpoint that queries user data and streams it as a CSV response.\\
\texttt{[Tool: edit\_file(path=``src/api/export.py'', changes=``create export\_csv endpoint with StreamingResponse'')]}\\
\texttt{[Tool: edit\_file(path=``src/api/routes.py'', changes=``register /api/export route'')]}\\[4pt]
\textbf{Turn 3 (User):} \textbf{[DRIFT]} Looks good. We're planning to offer this to enterprise customers next quarter---some of them have datasets with 10M+ rows.\\[4pt]
\textbf{Turn 4 (Agent):} Given the enterprise scale you mentioned, the current implementation would likely cause memory issues and timeouts with 10M+ row datasets. I'll add chunked streaming to avoid loading all data into memory, implement pagination with a cursor, add a progress callback mechanism, and include request timeout handling. I'll also add rate limiting to prevent abuse.\\
\texttt{[Tool: edit\_file(path=``src/api/export.py'', changes=``add chunked streaming with cursor-based pagination, configurable chunk\_size, timeout handling, rate limiting decorator'')]}\\
\texttt{[Tool: edit\_file(path=``tests/test\_export.py'', changes=``add test\_large\_dataset\_streaming and test\_export\_timeout'')]}\\
\texttt{[Tool: run\_tests(path=``tests/test\_export.py'')]}
\end{quote}

\noindent\textit{Analysis:} The developer never explicitly requests scalability improvements. However, the mention of enterprise customers and 10M+ row datasets implies requirements for chunked streaming, memory management, and timeout handling. The agent must proactively infer these needs and implement them without being asked.

\subsection{T6: Intent Backtracking}

\begin{quote}
\small
\textbf{Turn 1 (User):} Implement caching for the \texttt{get\_user\_profile()} function using Redis.\\
\textbf{Turn 1 (Agent):} I'll add Redis-based caching to the profile lookup.\\
\texttt{[Tool: read\_file(path=``src/services/user\_service.py'')]}\\
\texttt{[Tool: edit\_file(path=``src/services/user\_service.py'', changes=``add Redis cache layer with TTL=300s, cache-aside pattern'')]}\\
\texttt{[Tool: edit\_file(path=``requirements.txt'', changes=``add redis==5.0.0'')]}\\[4pt]
\textbf{Turn 2 (User):} Actually, let's use an in-memory LRU cache instead. Redis adds operational complexity we don't want right now.\\[4pt]
\textbf{Turn 2 (Agent):} I'll switch to \texttt{functools.lru\_cache} and remove the Redis dependency.\\
\texttt{[Tool: edit\_file(path=``src/services/user\_service.py'', changes=``replace Redis cache with @lru\_cache(maxsize=1024), remove redis import'')]}\\
\texttt{[Tool: edit\_file(path=``requirements.txt'', changes=``remove redis==5.0.0'')]}\\[4pt]
\textbf{Turn 3 (User):} \textbf{[DRIFT]} Hmm, on second thought, we do need distributed caching for our multi-instance deployment. Go back to the Redis approach you had initially.\\[4pt]
\textbf{Turn 3 (Agent):} I'll revert to the Redis-based implementation. Let me restore the cache-aside pattern with the TTL configuration from the first version.\\
\texttt{[Tool: edit\_file(path=``src/services/user\_service.py'', changes=``restore Redis cache-aside pattern with TTL=300s, remove lru\_cache'')]}\\
\texttt{[Tool: edit\_file(path=``requirements.txt'', changes=``re-add redis==5.0.0'')]}\\
\texttt{[Tool: run\_tests(path=``tests/test\_user\_service.py'')]}\\[4pt]
\textbf{Turn 4 (User):} Perfect. Also add cache invalidation when the profile is updated.\\
\textbf{Turn 4 (Agent):} I'll add invalidation logic to the \texttt{update\_user\_profile()} method.\\
\texttt{[Tool: edit\_file(path=``src/services/user\_service.py'', changes=``add cache.delete(user\_id) in update\_user\_profile'')]}\\
\texttt{[Tool: run\_tests(path=``tests/test\_user\_service.py'')]}
\end{quote}

\noindent\textit{Analysis:} The developer explicitly reverses a prior decision, requesting a return to the Redis approach after having switched to LRU cache. The agent must maintain a history of prior implementations and restore the Redis version accurately, including the dependency in \texttt{requirements.txt}. A subsequent refinement (T1: adding cache invalidation) tests whether the agent can seamlessly continue building on the restored state.


% ============================================================
\section{Repository Selection Criteria}
\label{app:repo-selection}

\subsection{Selection Criteria}

We select source repositories for \bench{} task construction according to five criteria:

\begin{enumerate}[leftmargin=*, itemsep=2pt]
    \item \textbf{Language:} Python, to align with the SWE-bench ecosystem and ensure compatibility with existing SWE agent tooling.
    \item \textbf{Popularity:} At least 1{,}000 GitHub stars, ensuring active maintenance, community familiarity, and sufficient code quality for meaningful evaluation.
    \item \textbf{Test coverage:} Comprehensive test suites with $>$60\% line coverage, enabling test-based verification of agent outputs. Repositories with minimal or flaky tests are excluded.
    \item \textbf{License:} Permissive open-source licenses (MIT, Apache-2.0, BSD) to ensure legal redistribution and modification of code artifacts within the benchmark.
    \item \textbf{Domain diversity:} Repositories are drawn from diverse application domains (web frameworks, data processing, CLI tools, scientific computing, DevOps utilities, ML libraries) to ensure that drift scenarios cover a representative range of SWE tasks.
\end{enumerate}

\subsection{Selected Repositories}

Table~\ref{tab:repos} lists the repositories included in \bench{} with key statistics.

\begin{table}[h]
\centering
\small
\caption{Repositories used in \bench{} task construction. Stars and LOC are approximate as of the benchmark construction date.}
\label{tab:repos}
\begin{tabular}{llrrrr}
\toprule
\textbf{Repository} & \textbf{Domain} & \textbf{Stars} & \textbf{Tests} & \textbf{LOC (k)} & \textbf{License} \\
\midrule
Flask         & Web framework      & 68k  & 520+  & 35  & BSD-3     \\
FastAPI       & Web framework      & 80k  & 900+  & 45  & MIT       \\
Pydantic      & Data validation    & 21k  & 1800+ & 55  & MIT       \\
Click         & CLI framework      & 16k  & 650+  & 20  & BSD-3     \\
Scikit-learn  & ML library         & 60k  & 2500+ & 280 & BSD-3     \\
Rich          & Terminal rendering & 50k  & 400+  & 30  & MIT       \\
\bottomrule
\end{tabular}
\end{table}

\noindent These six repositories span web development, data validation, CLI tooling, machine learning, and terminal UI---providing diverse SWE contexts in which intent drift can manifest. All repositories have active maintenance histories and well-documented contribution guidelines, ensuring that the code patterns used in our tasks are representative of production-quality Python software.


% ============================================================
\section{Task Construction Guidelines}
\label{app:task-construction}

\subsection{Seed Task Design}

Each seed task is designed to represent a realistic developer request grounded in the selected repository's actual codebase. Seed tasks are constructed by:

\begin{enumerate}[leftmargin=*, itemsep=2pt]
    \item \textbf{Issue mining:} Reviewing closed GitHub issues and pull requests from each repository to identify common task types (bug fixes, feature additions, refactoring, performance optimization, test improvements).
    \item \textbf{Task specification:} Writing a natural-language developer request that is sufficiently concrete to be actionable but sufficiently open-ended to admit multiple valid approaches. Each seed task specifies the target module(s), the desired outcome, and relevant context.
    \item \textbf{Feasibility validation:} Verifying that the task can be completed by an expert developer within a single session (approximately 30 minutes) and that the repository's test suite can validate the solution.
\end{enumerate}

\noindent We construct 20 seed tasks per repository, yielding 120 seeds in total.

\subsection{Drift Blueprint Specification}

Each seed task is paired with a drift blueprint that specifies:

\begin{itemize}[leftmargin=*, itemsep=2pt]
    \item \textbf{Drift type(s):} One or two drift types from the taxonomy (T1--T6).
    \item \textbf{Trigger condition:} The point in the interaction at which drift is injected, defined relative to the agent's progress (\eg, ``after the agent has edited the target file but before running tests'').
    \item \textbf{Drift content:} The specific requirement change the simulated developer introduces (\eg, ``add a constraint that the public API must not change'').
    \item \textbf{Expected adaptation:} A description of the ideal agent response, including which files should be modified, which changes should be rolled back, and whether clarification is appropriate.
    \item \textbf{Dimension annotations:} Values for all five drift dimensions (distance, explicitness, timing, frequency, reversibility).
\end{itemize}

\subsection{Test Case Construction}

For each drift scenario, we construct or extend the repository's test suite to verify correct drift handling:

\begin{enumerate}[leftmargin=*, itemsep=2pt]
    \item \textbf{Pre-drift tests:} Tests that verify the agent's work before drift is introduced. These establish a baseline of correct behavior.
    \item \textbf{Post-drift tests:} Tests that verify the final repository state after drift has been handled. These encode the developer's revised requirements.
    \item \textbf{Negative tests:} Tests that should \emph{fail} if the agent ignores the drift and continues with the original plan. These provide a discriminative signal between drift-aware and drift-unaware behavior.
\end{enumerate}

\noindent All tests are runnable via \texttt{pytest} and are validated by the automated verification layer (L1) of our pipeline.

\subsection{Inter-Annotator Agreement}

Two trained annotators independently review each completed task for correctness of drift placement, type labeling, and test adequacy. Table~\ref{tab:iaa-swe} reports agreement statistics.

\begin{table}[h]
\centering
\small
\caption{Inter-annotator agreement (Cohen's $\kappa$) for SWE task construction.}
\label{tab:iaa-swe}
\begin{tabular}{lc}
\toprule
\textbf{Annotation Dimension} & \textbf{Cohen's $\kappa$} \\
\midrule
Drift point placement & 0.84 \\
Drift type classification (T1--T6) & 0.79 \\
Test adequacy (pre/post/negative) & 0.77 \\
Expected adaptation correctness & 0.76 \\
Dimension annotations & 0.75 \\
\bottomrule
\end{tabular}
\end{table}

\noindent All agreement scores meet or exceed the $\kappa \geq 0.75$ threshold, indicating substantial inter-annotator reliability. Drift point placement shows the highest agreement ($\kappa = 0.84$), as the trigger conditions are precisely defined in the blueprint. Dimension annotations show the lowest agreement ($\kappa = 0.75$), reflecting occasional ambiguity in judging drift distance and reversibility for complex SWE tasks.


% ============================================================
\section{User Simulator Validation}
\label{app:user-simulator}

\subsection{Validation Protocol}

To ensure that the user simulator faithfully follows the drift blueprint without introducing artifacts that could confound evaluation, we plan a systematic validation study following the protocol established by $\tau$-Knowledge~\citep{tauknowledge2026}.

\paragraph{Trajectory sampling.}
We sample $N = 200$ complete trajectories (stratified by drift type, repository, and difficulty) from the user simulator interacting with a reference SWE agent. Each trajectory is recorded with full tool-call logs.

\paragraph{Manual annotation.}
Two expert annotators independently review each trajectory and classify simulator turns into three categories:

\begin{enumerate}[leftmargin=*, itemsep=2pt]
    \item \textbf{Correct:} The simulator's message faithfully follows the blueprint, introduces the drift at the specified trigger point, and provides a natural developer response.
    \item \textbf{Minor error:} The simulator's message is slightly unnatural or includes minor inconsistencies (\eg, referencing a file name not previously mentioned), but does not fundamentally alter the drift scenario.
    \item \textbf{Critical error:} The simulator introduces incorrect technical information, triggers drift at the wrong point, contradicts the blueprint's requirements, or produces a response that would mislead the agent in a way not intended by the scenario design.
\end{enumerate}

\paragraph{Target metrics.}
Following $\tau$-Knowledge, which achieved a critical error rate of 2\% across 500 trajectories, we target:

\begin{itemize}[leftmargin=*, itemsep=2pt]
    \item Critical error rate $<$ 5\% of all simulator turns.
    \item Minor error rate $<$ 15\% of all simulator turns.
    \item Blueprint adherence (fraction of drift events triggered at the correct point): $>$ 95\%.
\end{itemize}

\subsection{Preliminary Results}

In a pilot validation on 50 trajectories, the simulator achieves a critical error rate of 3.2\%, a minor error rate of 11.4\%, and blueprint adherence of 96.8\%. The most common critical errors involve the simulator providing overly specific technical hints that effectively solve the task for the agent, reducing the diagnostic value of the scenario. We address this by adding explicit instructions to the simulator prompt to avoid revealing implementation details beyond what the blueprint specifies. Full validation results will be reported upon completion of the 200-trajectory study.


% ============================================================
\section{Dataset License and Ethics}
\label{app:ethics}

\paragraph{Source repository licenses.}
All repositories used in \bench{} are open-source projects released under permissive licenses (MIT, BSD-3, Apache-2.0). Our benchmark uses these repositories' code structures and test patterns for evaluation purposes, which falls within the scope of permitted use under all applicable licenses. Table~\ref{tab:repos} lists the specific license for each repository.

\paragraph{No personal data.}
The benchmark contains no personally identifiable information. All developer names, usernames, commit messages, and issue descriptions appearing in task scenarios are synthetic. No real GitHub users, contributors, or issue reporters are referenced or represented. Code snippets used in task construction are either original or adapted from public repository code under permissive licenses.

\paragraph{Benchmark license.}
\bench{} is released under the \textbf{Creative Commons Attribution 4.0 International (CC-BY-4.0)} license. The benchmark dataset---including task descriptions, drift blueprints, test cases, and evaluation annotations---may be freely shared and adapted for any purpose, provided appropriate credit is given to the original authors.

\paragraph{Code license.}
The evaluation framework, user simulator, and analysis scripts are released under the \textbf{MIT License}. This includes the blueprint-to-trajectory pipeline, the automated verification tools, and the metric computation code.

\paragraph{Intended use and limitations.}
The benchmark is intended exclusively for research evaluation of SWE agents under intent drift. It should not be used as training data for production systems without careful consideration of the synthetic nature of the drift scenarios. The test cases are designed for diagnostic evaluation and may not reflect the full distribution of real-world developer requirements.

\paragraph{Potential risks.}
We identify two primary risks. First, \emph{benchmark overfitting}: models may be optimized to perform well on \bench{} patterns without developing genuine drift-handling capabilities. We mitigate this by releasing the generation pipeline alongside the static dataset, enabling researchers to produce fresh scenarios. Second, \emph{synthetic distribution gap}: despite three-layer verification, synthetic drift scenarios may not perfectly reflect organic developer behavior. We encourage complementary evaluation with real user studies and plan to release updated benchmark versions as validation data accumulates.

% ============================================================

\end{document}
